 &\le2B_{\rm cm}+4096(1028+5D+84000D^2)=C_{\rm cm}.
\end{align*}
For the last line use $\log28<4$, $\log d\le D$, and
$\log(2B_{\rm cm})\le2+2D+42000D^2$.
No bound for the degree of the auxiliary good-reduction extension
enters the isogeny estimate.

\paragraph{Non-CM fibers and scalar descent.}
Now suppose $\operatorname{End}^0(A_{\overline\Q})=B$.
In particular $A$ is absolutely simple. For each
$\sigma\in\operatorname{Gal}(L/\Q)$ choose a geometric isomorphism
$\psi_\sigma:{}^\sigma A\to A$. The automorphism
$b\mapsto\psi_\sigma{}^\sigma b\psi_\sigma^{-1}$ of $B$ is inner.
Correcting it by Skolem--Noether and clearing denominators gives
a $B$-linear isogeny $\mu_\sigma:{}^\sigma A\to A$, defined over
$L$ because all the relevant Hom groups are defined there. Take
$\mu_1=1$ and inflate this finite choice to $G_\Q$.
The discrepancy
\[
 c(\sigma,\tau)=
 \mu_\sigma{}^\sigma\mu_\tau\mu_{\sigma\tau}^{-1}
\]
commutes with $B$, hence belongs to $\Q^\times$.
Set $a_\sigma=(\deg\mu_\sigma)^{-1/4}>0$ and
$\nu_\sigma=a_\sigma\mu_\sigma$ in the category with algebraic
coefficients. Throughout these coefficient extensions Galois acts on
the varieties and morphisms and acts trivially on the added scalars;
the induced homology maps are extended linearly. In particular, the
$a_\sigma$ are constant coefficients in the products below.
Multiplicativity of degree and
$\deg[c]=c^4$ show that
\[
 \nu_\sigma{}^\sigma\nu_\tau
       =\varepsilon(\sigma,\tau)\nu_{\sigma\tau},
 \qquad \varepsilon(\sigma,\tau)\in\{1,-1\}.
\]
Thus $\varepsilon$ is a normalized sign cocycle inflated from
$\operatorname{Gal}(L/\Q)$. Only finitely many algebraic
coefficients $a_\sigma$ occur.

We need a splitting with controlled ramification, not merely
existence of a projective lift. Identify
$H^2(G_\Q,\{1,-1\})$ with $\operatorname{Br}(\Q)[2]$.
At a finite prime outside $T$ the class of $\varepsilon$ is zero:
its restriction is unramified and the unramified local quotient
has cohomological dimension one. Let $T_-$ be the finite primes
where its invariant is $1/2$. For each odd $p\in T_-$ choose an
odd character of $(\Z/p\Z)^\times$; at $2\in T_-$ use the
nontrivial character modulo four. Their product is a Dirichlet
character $\chi$, unramified outside $T$, whose order $m$ satisfies
\[
 m\le2\prod_{\substack{p\in T\\p\ne2}}(p-1)
       \le2R^R M.
\]
Here the product over an empty set is one.

To check that $\chi$ has the required square-root obstruction,
use the exact sequence of constant discrete Galois modules
$1\to\{\pm1\}\to\mu_\infty\xrightarrow{z\mapsto z^2}
\mu_\infty\to1$. By local reciprocity, a finite-order character
of $G_{\Q_p}$ has a character square root exactly when its value
on $-1\in\Q_p^\times$ is $1$. Indeed this is the only obstruction
on the torsion part of $\Q_p^\times$; the valuation factor and
the procyclic principal-unit factors admit such roots. Hence the
boundary of $\chi$ and $[\varepsilon]$ agree at every finite
place. Brauer reciprocity makes them agree at the real place,
and the injective global-to-local Brauer map makes them equal
globally \cite[Chapter I]{milne}.
Choosing a continuous pointwise square root of $\chi$ and
correcting it by a sign cochain consequently gives a continuous
cochain $\beta:G_\Q\to\mu_{2m}$ such that
\[
 \beta^2=\chi,\qquad
 \beta(\sigma)\beta(\tau)\beta(\sigma\tau)^{-1}
        =\varepsilon(\sigma,\tau).
\]
Its restriction to $G_L$ is a character invariant under $G_\Q$.
Outside $T$, its inertia restrictions are quadratic and nontrivial
at only finitely many primes. Twisting by the quadratic character
of the product of those odd primes kills all these restrictions:
at an odd prime there is a unique nontrivial quadratic character
on inertia. This introduces ramification only at those primes
and possibly at $2\in T$, and changes neither displayed identity.
We henceforth use this choice of $\beta$.

Let $L'$ be the kernel field of $\beta|_{G_L}$.
Conjugacy invariance shows that $L'/\Q$ is Galois, and
\[
 [L':\Q]\le2md\le4DR^R M\le QM.
\]
The last inequality follows from $R\ge\max(D,10^{12})$ and
